% !TEX root = ../Projektdokumentation.tex
\section{Kundendokumentation}
\label{sec:Kundendokumentation}
Die Anwendung wird über einen Webbrowser aufgerufen. Nach erfolgreicher
Anmeldung stehen die Menüpunkte Dashboard, Organisationen, Benutzer,
Mail-Adressen, IP-Adressen und Audit-Log zur Verfügung. Welche Bereiche
sichtbar und bearbeitbar sind, richtet sich nach der zugewiesenen Rolle.

Neue Organisationen, Benutzer sowie Mail- und IP-Einträge werden über die
jeweiligen Listenansichten angelegt. Änderungen erfolgen über Bearbeitungs-
formulare, die Eingaben bereits beim Speichern validieren. Fehlerhafte
Eingaben werden direkt am betroffenen Feld angezeigt. Dadurch kann die Pflege
ohne technische Detailkenntnisse nachvollziehbar und sicher erfolgen.

Das Audit-Log dient der Rückverfolgung fachlicher Änderungen. Es zeigt, wann
ein Datensatz erstellt, geändert oder gelöscht wurde. Soweit fachlich
freigegeben, können Änderungen aus dem Audit-Log heraus zurückgesetzt
werden. Ein kurzer Benutzerhinweis mit den wichtigsten Bedienregeln befindet
sich zusätzlich im \Anhang{app:Kundenhinweis}.

Für die Einführung beim Kunden ist keine umfangreiche Schulung erforderlich.
Die Anwendung orientiert sich an den bekannten Objekten des Fachprozesses und
verwendet eine durchgängige Bedienlogik für Listen-, Detail- und
Bearbeitungsansichten. Dadurch kann die Kundendokumentation auf die für den
Betrieb notwendigen Schritte konzentriert bleiben.

Im laufenden Betrieb ist darauf zu achten, dass Berechtigungen nur nach dem
tatsächlichen Aufgabenbereich vergeben werden. Insbesondere administrative
Rollen dürfen nicht pauschal verteilt werden, da sie organisationsübergreifende
Pflegevorgänge ermöglichen. Neue Einträge für Mail-Adressen und IP-Adressen
sollen vor dem Speichern fachlich geprüft werden, damit nur berechtigte
Ausnahmen in den verwalteten Datenbestand aufgenommen werden.
